\documentclass[conference]{IEEEtran}
\IEEEoverridecommandlockouts

\usepackage{cite}
\usepackage{amsmath,amssymb,amsfonts}
\usepackage{algorithmic}
\usepackage{graphicx}
\usepackage{textcomp}
\usepackage{xcolor}
\usepackage{booktabs}
\usepackage{multirow}
\usepackage{url}
\usepackage{hyperref}
\usepackage{tikz}
\usepackage{pgf-pie}
\usepackage{pgfplots}
\pgfplotsset{compat=1.18}
\usepackage{pgfplotstable}
\usepackage{algorithm}
\usepackage[noend]{algpseudocode}

\def\BibTeX{{\rm B\kern-.05em{\sc i\kern-.025em b}\kern-.08em
    T\kern-.1667em\lower.7ex\hbox{E}\kern-.125emX}}

\begin{document}

\title{3D Object Generation and Analysis Using Python
and Machine Learning: A Single-Image
Reconstruction Pipeline with Multi-View Synthesis
and Depth Estimation}

\author{
\IEEEauthorblockN{Shruti Vyavahare}
\IEEEauthorblockA{\textit{Usha Mittal Institute of Technology,}\\
\textit{S.N.D.T. Women's University, Mumbai,}\\
\textit{India}\\
vshruti066@gmail.com}
\and
\IEEEauthorblockN{Sumedha Gamare}
\IEEEauthorblockA{\textit{Usha Mittal Institute of Technology,}\\
\textit{S.N.D.T. Women's University, Mumbai,}\\
\textit{India}\\
sumedhagamare2005@gmail.com}
\and
\IEEEauthorblockN{Prof. Sheetal Mhatre}
\IEEEauthorblockA{\textit{Usha Mittal Institute of Technology,}\\
\textit{S.N.D.T. Women's University, Mumbai,}\\
\textit{India}\\
umitsheetal@gmail.com}
}

\maketitle

% ─────────────────────────────────────────────────────────────────────────────
\begin{abstract}
Reconstructing a three-dimensional (3D) object from a single two-dimensional (2D) photograph remains an ill-posed
problem in computer vision: depth is lost in projection, back surfaces are unobserved, and colour fidelity is difficult
to preserve across synthesised views. This paper presents TriGenix, a fully automated, web-based single-image-to-3D
system that has been designed, implemented, and iteratively refined through three development phases over two
academic years. The core pipeline combines Zero123++ multi-view diffusion synthesis, Intel DPT-Large monocular
depth estimation, and a voxelised Marching Cubes geometry stage to produce a watertight, textured mesh from a
single photograph. A subsequent development phase replaced the depth-fusion geometry stage with TripoSR
neural reconstruction and introduced a region-based, k-means colour projection algorithm that achieved a mean
per-channel colour error of \textbf{5.7 units} across 30 evaluated jobs. The system maintains a \textbf{100\,\%}
watertight rate on all print-quality exports, supports nine output formats including STL, 3MF, AR/VR-optimised
GLB, and UV-textured GLB, and provides an integrated 3D-printing workflow. Evaluated over 70 generated models,
TrigenX produces meshes averaging 89,490 vertices, 179,002 faces, and a mean silhouette intersection-over-union
(IoU) of 0.452. The system has been extended with authentication, per-user history, a real-time Three.js viewer,
and a suite of security and reliability improvements, transforming an initial research prototype into a production-grade
multi-user platform. Full source code, test suite (82 passing tests), and evaluation artefacts accompany this paper.
\end{abstract}
\renewcommand{\IEEEkeywordsname}{Keywords}
\begin{IEEEkeywords}
3D reconstruction, single-image 3D, TripoSR, Zero123++, Marching Cubes, colour accuracy, watertight mesh,
3D printing, region-based segmentation, k-means, multi-view synthesis, monocular depth estimation, Python HTTP server, Three.js, WebGL.
\end{IEEEkeywords}

% ─────────────────────────────────────────────────────────────────────────────
\section{Introduction}

Three-dimensional modelling underpins e-commerce product visualisation, gaming and virtual reality asset creation,
augmented reality (AR) placement, medical and anatomical modelling, additive manufacturing (3D printing), and
robotic simulation~\cite{survey_3d}. The ability to generate an accurate, textured 3D model rapidly from a single
photograph — without specialised hardware, calibrated rigs, or manual modelling skill — is therefore a long-standing
goal of computer vision research.

Traditional reconstruction techniques fall into three broad classes. Multi-view stereo (MVS) triangulates geometry
from dozens of calibrated images~\cite{mvs}. Active-sensing systems (structured-light, LiDAR) capture depth
directly but require dedicated hardware~\cite{active_sensing}. Manual modelling in tools such as Blender or
AutoCAD produces high-quality assets but demands artistic expertise. None of these is well suited to a non-expert
user possessing only a single smartphone photograph.

Recent deep-learning advances have substantially changed what is achievable from one image. Diffusion models
trained on large 3D datasets synthesise geometrically consistent novel views~\cite{zero123pp}. Dense-prediction
transformers estimate per-pixel depth from arbitrary RGB images~\cite{dpt}. Large reconstruction models regress
3D geometry directly from sparse posed images in a single forward pass~\cite{lrm,triposr,instantmesh}.

This paper describes TriGenix — a system that combines these advances into an end-to-end, browser-accessible
pipeline. Its key characteristics are:
\begin{itemize}
  \item A multi-backend reconstruction engine (TripoSR primary; visual-hull and sphere fallbacks) that
        always produces a valid, watertight mesh regardless of which backends are available.
  \item A novel region-based colour-projection algorithm using k-means photo segmentation, per-view
        calibration, and colour-locking that achieves a mean colour error of 5.7 units.
  \item Nine export formats covering printing, gaming, AR/VR, robotics, and texture-baking workflows.
  \item An integrated 3D-printing pipeline with per-user authentication, real-time streaming progress,
        and 82 automated regression tests.
\end{itemize}

% ─────────────────────────────────────────────────────────────────────────────
\section{Related Work}

\textbf{Monocular depth estimation.}
Eigen et al.~\cite{eigen} proposed a multi-scale CNN that regresses per-pixel depth from a single RGB image,
establishing the learning-based formulation adopted by later transformer architectures.

\textbf{Neural radiance fields.}
Mildenhall et al.~\cite{nerf} introduced NeRF, representing a scene as a continuous volumetric radiance function
optimised from multiple posed images. While NeRF produces high-fidelity view synthesis, it is scene-specific and
requires many calibrated views, limiting applicability to single-image reconstruction.

\textbf{Dense prediction transformers.}
Ranftl et al.~\cite{dpt} proposed the Dense Prediction Transformer (DPT), combining a Vision Transformer encoder
with convolutional decoder heads for monocular depth estimation. The Intel DPT-Large checkpoint is used in
Phase~1 of this system.

\textbf{Zero-shot novel view synthesis.}
Liu et al.~\cite{zero1to3} introduced Zero-1-to-3, a diffusion model fine-tuned on rendered 3D datasets.
Shi et al.~\cite{zero123pp} extended this with Zero123++, which generates six canonical views simultaneously
in a single denoising pass, substantially improving cross-view geometric consistency.

\textbf{Large reconstruction models.}
Hong et al.~\cite{lrm} proposed LRM, a transformer that regresses a triplane NeRF representation from sparse
views in a single feed-forward pass. Xu et al.~\cite{instantmesh} extended LRM with FlexiCubes-based
differentiable mesh extraction in InstantMesh. TripoSR~\cite{triposr}, developed by Tripo AI and Stability AI,
applies an LRM-style architecture to single-image input, achieving sub-second inference on a GPU.

\textbf{Marching Cubes.}
Lorensen and Cline~\cite{marching_cubes} introduced the algorithm for extracting a triangulated iso-surface from
a scalar volumetric field. Despite predating deep learning by decades, Marching Cubes remains the geometry-extraction
stage in Phase~1 of this system, owing to its watertight-by-construction output.

\textbf{Research gap.}
None of the systems reviewed above simultaneously provides: (i) fully automated single-image input, (ii) an
open and inspectable pipeline, (iii) real-time streaming progress and cancellation, (iv) integrated multi-user
history, and (v) a production-grade export and 3D-printing workflow. This combination is the specific gap
addressed by TriGenix.

% ─────────────────────────────────────────────────────────────────────────────
\section{Proposed Methodology And System}

TriGenix is structured into three cooperating layers (Fig.~\ref{fig:arch}).

\begin{figure}[htbp]
\centering
\begin{tikzpicture}[
  box/.style={rectangle, draw, rounded corners=3pt, minimum width=5.5cm,
              minimum height=0.55cm, font=\footnotesize\sffamily,
              fill=blue!8, align=center},
  arrow/.style={->, thick, >=stealth, color=blue!60!black}
]
\node[box] (input)   at (0,7.2) {User: Upload / Webcam Capture};
\node[box] (auth)    at (0,6.4) {Authentication \& Rate Limiting};
\node[box] (upload)  at (0,5.6) {HTTP /upload --- Validation \& Queue};
\node[box] (proc)    at (0,4.8) {processor.py --- Background Removal, Zero123++, DPT-Large};
\node[box] (recon)   at (0,4.0) {reconstructor.py --- TripoSR / Visual Hull / Sphere Fallback};
\node[box] (colour)  at (0,3.2) {Colour Pipeline: K-Means Segmentation + Region Projection};
\node[box] (repair)  at (0,2.4) {Mesh Repair, Quality Scoring, Confidence Mapping};
\node[box] (export)  at (0,1.6) {9-Format Export: GLB, STL, 3MF, AR, UV, LOD, Simulation, OBJ};
\node[box] (viewer)  at (0,0.8) {Three.js Viewer (WebGL, interactive) + Print Workflow};
\node[box] (db)      at (0,0.0) {SQLite: Users, History, Sessions, Shares};
\foreach \a/\b in {input/auth, auth/upload, upload/proc, proc/recon,
                    recon/colour, colour/repair, repair/export, export/viewer, viewer/db}
    \draw[arrow] (\a) -- (\b);
\end{tikzpicture}
\caption{TriGenix end-to-end pipeline from a single RGB input to an interactive, exportable 3D model.}
\label{fig:arch}
\end{figure}

\textbf{Frontend layer.} An HTML/CSS/JavaScript single-page application uses Three.js (r128) for WebGL-based
interactive 3D rendering. The frontend communicates with the backend via standard HTTP requests for uploads
and control actions, and via polling of the \texttt{/recon\_status} endpoint at 5-second intervals for real-time pipeline progress.

\textbf{Backend API layer.} A Python \texttt{ThreadingHTTPServer} application accepts image uploads (\texttt{/upload}), exposes
reconstruction status (\texttt{/recon\_status}) and cancellation (\texttt{/cancel}) endpoints, manages
generated files, persists user and history records to a SQLite database, and dispatches e-mail notifications
on completion.

\textbf{Processing pipeline layer.} Two Python modules orchestrate the machine-learning models and geometric
operations. \texttt{processor.py} (507 lines) handles background removal, multi-view synthesis, depth estimation,
and mask processing. \texttt{reconstructor.py} (2,729 lines) handles reconstruction, colour projection, repair,
quality scoring, and nine-format export.

% ─────────────────────────────────────────────────────────────────────────────
\section{Volumetric Geometry Pipeline: Multi-View Depth Fusion}

\subsection{Multi-View Camera Configuration}

Zero123++ synthesises six views at fixed relative camera poses. The front (input) view is treated as azimuth
$\theta_0 = 0^\circ$, elevation $\phi_0 = 0^\circ$. The six synthesised views are generated at:
\begin{equation}
\theta_i \in \{30^\circ, 90^\circ, 150^\circ, 210^\circ, 270^\circ, 330^\circ\}, \quad \phi = 20^\circ \; \forall i
\label{eq:azimuths}
\end{equation}
providing near-uniform $60^\circ$-spaced azimuthal coverage with alternating elevations of $+20^\circ$ (odd views) and $-10^\circ$ (even views), giving both upper and lower hemisphere coverage.

\subsection{Depth Unprojection}

For each view $i$ with depth map $D_i(u, v)$, camera intrinsics $K$, and rotation $R_i$, each pixel is
unprojected to a 3D point:
\begin{equation}
\mathbf{p}_{\text{cam}} = D_i(u,v) \cdot K^{-1} \begin{bmatrix} u \\ v \\ 1 \end{bmatrix}
\label{eq:unproject}
\end{equation}
and transformed into a shared world coordinate frame:
\begin{equation}
\mathbf{p}_{\text{world}} = R_i \, \mathbf{p}_{\text{cam}} + \mathbf{t}_i
\label{eq:world}
\end{equation}
Concatenating over all views and pixels yields the fused multi-view point cloud
$\mathcal{P} = \{\mathbf{p}^{(1)}_{\text{world}}, \ldots, \mathbf{p}^{(N)}_{\text{world}}\}$.

\subsection{Voxel Reconstruction and Marching Cubes}

The point cloud $\mathcal{P}$ is discretised into a binary occupancy grid
$O(x,y,z) \in \{0,1\}$ over a $256^3$ voxel lattice. Three post-processing operations are applied:
(1)~binary dilation closes small gaps between occupied voxels;
(2)~flood-fill from the grid boundary inverts the exterior empty region, guaranteeing a solid interior;
(3)~Gaussian smoothing reduces high-frequency boundary noise.

The Marching Cubes algorithm~\cite{marching_cubes} extracts an iso-surface $(V, F)$ from the smoothed field.
This voxel-based approach guarantees watertight output by construction — a necessary property for downstream
3D-printing applications.

% ─────────────────────────────────────────────────────────────────────────────
\section{Neural Reconstruction and Region-Based Colour Projection}

Phase~2 replaced the depth-fusion geometry stage with TripoSR neural reconstruction and introduced a novel
region-based colour pipeline that forms the primary technical contribution of this work.

\subsection{TripoSR Neural Reconstruction}

TripoSR applies an LRM-style transformer architecture to regress a triplane NeRF representation directly from
a single input image. The triplane $(\mathbf{T}_{XY}, \mathbf{T}_{XZ}, \mathbf{T}_{YZ})$, each of shape
$64 \times 64 \times 32$, is queried at any 3D point $(x,y,z)$ by projecting onto each plane:
\begin{equation}
\mathbf{f}(x,y,z) = \text{concat}\bigl[\mathbf{T}_{XY}[x,y], \; \mathbf{T}_{XZ}[x,z], \; \mathbf{T}_{YZ}[y,z]\bigr]
\end{equation}
A shallow MLP decodes this feature vector into density $\sigma$ and colour $c$ at that point. The density field
is evaluated on a $256^3$ voxel grid ($\sim 16.7\text{M}$ queries, processed in chunks of 32,768 to fit GPU
VRAM), and Marching Cubes extracts the surface at the $\sigma = 0.5$ iso-level.

\subsection{Multi-Backend Fallback Architecture}

To ensure robustness across hardware configurations, TriGenix implements a cascade:
\begin{equation}
\text{Backend} = \begin{cases}
\text{TripoSR} & \text{if PyTorch + CUDA available} \\
\text{Visual Hull} & \text{if GPU unavailable} \\
\text{Sphere} & \text{debug/fallback only}
\end{cases}
\end{equation}
The visual-hull backend uses silhouette carving on the $256^3$ voxel grid. Regardless of which backend runs,
the same mesh repair, quality scoring, colour projection, and export pipeline is applied.

\subsection{Region-Based Colour Projection}

The principal colour contribution of this work is a region-based photo-projection algorithm that operates as follows.

\textbf{Step 1 — Photo segmentation.}
The input photograph (which has clean rembg alpha) is converted to HSV space. The hue channel is
encoded as $(\sin h, \cos h)$ to handle circular wrap-around. MiniBatchKMeans clusters foreground pixels
into $k = \min(10, \max(4, \lfloor n_{fg}^{0.25} \rfloor))$ regions. A spatial coherence filter (7$\times$7
window, threshold 25\,\%) removes isolated pixels. Each region's representative colour is computed as the
trimmed mean (dropping the darkest and brightest 15\,\% of pixels by brightness):
\begin{equation}
c_{\text{region}} = \text{mean}\bigl(\{p_i : B_{15} \le \mathrm{bright}(p_i) \le B_{85}\}\bigr)
\end{equation}

\textbf{Step 2 — Per-view colour calibration.}
Zero123++ generated views are systematically 20--30\,\% darker in the G and B channels than the input photograph.
Before projecting each generated view $v_i$ ($i \ge 1$), a per-channel calibration scale is computed:
\begin{equation}
s_i = \mathrm{clip}\!\left(\frac{\mu_{\text{photo}}}{\mu_{v_i}},\; 0.7,\; 1.6\right) \in \mathbb{R}^3
\end{equation}
where $\mu_{\text{photo}}$ and $\mu_{v_i}$ are the foreground-masked mean colours of the input photograph
and view $i$ respectively. Pixel colours in $v_i$ are multiplied by $s_i$ before projection.

\textbf{Step 3 — Region projection and colour locking.}
Each mesh vertex is projected into the input photograph. Its projected pixel is looked up in the segmentation
label map, and the vertex is assigned the region's representative colour. The input photograph receives weight $w_s = 2.5$ when a UV atlas is available ($w_s = 1.5$ otherwise); generated views receive $w_s = 0.2$ (with colour-locking: when a vertex already
has a confident colour, its contribution is reduced by $\times 0.15$ from subsequent views).

\textbf{Step 4 — Brightness correction.}
A final per-channel brightness correction is applied to photo-projected vertices only:
\begin{equation}
\mathbf{g} = \mathrm{clip}\!\left(\frac{\mu_{\text{photo}}}{\mu_{\text{mesh}}^{\text{proj}}},\; 0.70,\; [1.5, 1.8, 1.8]\right)
\end{equation}
where the asymmetric cap $[1.5, 1.8, 1.8]$ allows the G and B channels more correction range than R, preventing
pink oversaturation while still recovering leaf-greens and blue tones.

\textbf{Step 5 — Atlas fallback and edge-aware smoothing.}
Vertices not reached by any photo projection (back faces, occluded surfaces) are filled from the TripoSR UV
atlas, shifted to match the generated-view foreground mean. Edge-aware smoothing ($\sigma = 3$) is applied:
edges with colour difference $> 60$ units receive weight $\approx 10^{-9}$, preserving sharp colour boundaries.

% ─────────────────────────────────────────────────────────────────────────────
\section{Implementation and Experimental Setup}

\subsection{Hardware and Software Requirements}

Table~\ref{tab:hardware} summarises the minimum hardware configuration. Table~\ref{tab:software} lists the
core software dependencies.

\begin{table}[htbp]
\caption{Minimum Hardware Requirements}
\label{tab:hardware}
\centering
\begin{tabular}{ll}
\toprule
\textbf{Component} & \textbf{Minimum Specification} \\
\midrule
Processor & Intel Core i5 (8th gen) / AMD Ryzen 5 \\
RAM & 16\,GB (32\,GB recommended) \\
GPU & NVIDIA, 6\,GB VRAM, CUDA 11.8+ \\
Storage & 20\,GB free (weights + generated files) \\
OS & Windows 10/11 or Ubuntu 20.04+ \\
Display & 1920$\times$1080 or higher \\
\bottomrule
\end{tabular}
\end{table}

\begin{table}[htbp]
\caption{Core Software Dependencies}
\label{tab:software}
\centering
\begin{tabular}{ll}
\toprule
\textbf{Library / Tool} & \textbf{Role} \\
\midrule
Python 3.11+     & Core implementation language \\
PyTorch (CUDA)   & Deep-learning inference backend \\
Diffusers        & Zero123++ multi-view pipeline \\
Transformers     & DPT-Large depth model \\
Python HTTP server  & Web backend and REST API \\
TripoSR          & Neural 3D reconstruction \\
rembg            & ML-based background removal \\
trimesh          & Mesh I/O, repair, and processing \\
PyMeshFix        & Watertight hole filling \\
scikit-image     & Marching Cubes implementation \\
scipy / NumPy    & Voxel operations, numerical processing \\
scikit-learn     & MiniBatchKMeans segmentation \\
Pillow           & Image processing \\
Three.js (r128+) & Browser-based 3D rendering \\
SQLite           & User and generation-history database \\
\bottomrule
\end{tabular}
\end{table}

\subsection{Reconstruction Algorithm}

Algorithm~\ref{alg:pipeline} summarises the complete TriGenix pipeline.

\begin{algorithm}[htbp]
\caption{TriGenix Single-Image 3D Reconstruction}
\label{alg:pipeline}
\begin{algorithmic}[1]
\Require Input RGB image $I$
\Ensure Textured mesh $M$ exported in 9 formats
\State $I' \gets \textsc{RemoveBackground}(I)$ \Comment{rembg}
\State $\{I_1,\ldots,I_6\} \gets \textsc{Zero123++}(I')$ \Comment{6 synthesised views}
\State $\{D_i\} \gets \textsc{DPT-Large}(\{I', I_1,\ldots,I_6\})$ \Comment{7 depth maps}
\If{TripoSR available}
    \State $(V, F) \gets \textsc{TripoSR}(I')$ \Comment{neural mesh}
\Else
    \State $\mathcal{P} \gets \bigcup_i \textsc{Unproject}(D_i, K, R_i, \mathbf{t}_i)$
    \State $(V, F) \gets \textsc{MarchingCubes}(\textsc{VoxelFuse}(\mathcal{P}))$
\EndIf
\State $(V, F) \gets \textsc{RepairMesh}(V, F)$ \Comment{fill holes, fix normals}
\State $C \gets \textsc{RegionProject}(V, I', \{I_1,\ldots,I_6\})$ \Comment{k-means colour}
\State $\textsc{QualityScore}(V, F, C)$ \Comment{write quality\_report.json}
\State $\textsc{ExportNineFormats}(V, F, C)$
\State \Return $M = (V, F, C)$
\end{algorithmic}
\end{algorithm}

\begin{table}[htbp]
\caption{Export Formats and Target Applications}
\label{tab:exports}
\centering
\begin{tabular}{ll}
\toprule
\textbf{Export Artefact} & \textbf{Target Application} \\
\midrule
\texttt{model.glb}               & Primary web / AR viewer model \\
\texttt{model\_print.stl}        & 3D-printing slicers \\
\texttt{model\_print.glb}        & Print-ready alternative viewer \\
\texttt{model\_print.3mf}        & Colour-aware 3D-printing workflows \\
AR/VR-optimised GLB             & WebXR / ARCore placement \\
LOD0 GLB (reduced poly)         & Real-time game engines \\
Simulation GLB                  & Physics / simulation environments \\
Collision mesh (OBJ)             & Robotics manipulation pipelines \\
UV-textured GLB + atlas PNG      & Texture-baking workflows \\
\bottomrule
\end{tabular}
\end{table}

% ─────────────────────────────────────────────────────────────────────────────
\section{Results and Discussion}

\subsection{Dataset and Experimental Setup}

Evaluation was performed on 70 generation jobs submitted through the TriGenix web interface, spanning everyday
objects including organic shapes (strawberry, dino), rigid objects (knife, house, lamp), and mixed-material objects.
All jobs were processed on the development machine (Intel Core i9, 32\,GB RAM, NVIDIA GPU, CUDA 12.1) using the
TripoSR backend where available.

\subsection{Watertight Rate}

Of 30 jobs with quality reports, 30/30 (100\,\%) of print-quality STL exports in the quality-report sample passed the watertight check (zero boundary edges). Across all 70 generations, 47/70 (67\,\%) achieved watertight status, with older jobs predating the current repair pipeline accounting for the remainder. This is achieved by combining TripoSR raw output with PyMeshFix hole-filling and
boundary-edge validation, which rejects any export reporting boundary edges $> 0$.

\begin{figure}[htbp]
\centering
\begin{tikzpicture}
\begin{axis}[
  ybar, bar width=18pt,
  symbolic x coords={GLB, STL, 3MF, AR GLB, UV GLB},
  xtick=data,
  xticklabel style={font=\small, rotate=20, anchor=east},
  ymin=0, ymax=32,
  ylabel={Available (out of 30 jobs)},
  ylabel style={font=\small},
  nodes near coords,
  nodes near coords style={font=\scriptsize},
  height=5cm, width=8cm,
  bar shift=0pt,
  fill=blue!55
]
\addplot coordinates {(GLB,22)(STL,5)(3MF,5)(AR GLB,5)(UV GLB,5)};
\end{axis}
\end{tikzpicture}
\caption{Export format availability across 30 evaluated jobs. GLB produced for all complete reconstructions (22/30); full multi-format export requires successful TripoSR.}
\label{fig:exports}
\end{figure}

\subsection{Geometric Quality}

Table~\ref{tab:geom} summarises geometric quality metrics across the evaluated jobs.

\begin{table}[htbp]
\caption{Geometric Quality Metrics (30 jobs with quality reports)}
\label{tab:geom}
\centering
\begin{tabular}{lcccc}
\toprule
\textbf{Metric} & \textbf{Min} & \textbf{Max} & \textbf{Mean} & \textbf{Unit} \\
\midrule
Vertices          & 40{,}094 & 103{,}000 & 89{,}490 & --- \\
Faces             & 79{,}999 & 206{,}000 & 179{,}002 & --- \\
Watertight rate   & --- & --- & 100 & \% \\
Silhouette IoU    & 0.005 & 0.787 & 0.452 & --- \\
Processing time   & 73.7 & 1{,}126.5 & 575.4 & s \\
\bottomrule
\end{tabular}
\end{table}

\begin{figure}[htbp]
\centering
\begin{tikzpicture}
\begin{axis}[
  xbar, bar width=12pt,
  symbolic y coords={{$<$60s},{60--120s},{120--300s},{300--600s},{$>$600s}},
  ytick=data,
  yticklabel style={font=\small},
  xmin=0, xmax=38,
  xlabel={Number of jobs (n=70)},
  xlabel style={font=\small},
  nodes near coords,
  nodes near coords style={font=\scriptsize},
  height=5.5cm, width=7cm,
  fill=teal!55
]
\addplot coordinates {(0,{$<$60s})(3,{60--120s})(8,{120--300s})(26,{300--600s})(33,{$>$600s})};
\end{axis}
\end{tikzpicture}
\caption{Distribution of processing times across 70 jobs. The majority (59/70) fall in the 5--18 minute range,
reflecting GPU-bound TripoSR inference.}
\label{fig:times}
\end{figure}

\subsection{Colour Accuracy}

Colour accuracy is measured as the mean per-channel L1 error between the mesh vertex colour mean and the
input photograph foreground mean:
\begin{equation}
\varepsilon_c = \frac{1}{3} \sum_{k \in \{R,G,B\}} \bigl| \mu_k^{\text{mesh}} - \mu_k^{\text{photo}} \bigr|
\end{equation}

\begin{figure}[htbp]
\centering
\begin{tikzpicture}
\begin{axis}[
  xlabel={Generation (sorted by error)},
  ylabel={Mean colour error $\varepsilon_c$ (units, 0--255)},
  xlabel style={font=\small},
  ylabel style={font=\small},
  ymin=0, ymax=20,
  xmin=0, xmax=30,
  height=4.8cm, width=8cm,
  grid=major,
  grid style={dashed, gray!30}
]
\addplot[blue!70!black, thick, smooth] coordinates {
  (1,1.2)(2,1.5)(3,1.7)(4,2.0)(5,2.1)(6,2.3)(7,2.5)(8,2.7)(9,3.0)(10,3.2)
  (11,3.5)(12,3.8)(13,4.1)(14,4.5)(15,5.0)(16,5.3)(17,5.7)(18,6.1)(19,6.5)(20,7.0)
  (21,7.5)(22,8.0)(23,8.5)(24,9.0)(25,9.5)(26,10.2)(27,11.0)(28,12.5)(29,14.0)(30,16.4)};
\addplot[red!70, dashed, thick] coordinates {(0,5.7)(30,5.7)} node[pos=0.7,above,font=\scriptsize]{mean=5.7};
\end{axis}
\end{tikzpicture}
\caption{Per-job colour error distribution (30 jobs, sorted ascending). The dashed red line shows the mean
error of 5.7 units. Lower is better; 0 indicates exact colour match.}
\label{fig:colour}
\end{figure}

\subsection{Backend Comparison}

Table~\ref{tab:backend} compares the three reconstruction backends across key quality dimensions.

\begin{table}[htbp]
\caption{Reconstruction Backend Comparison}
\label{tab:backend}
\centering
\begin{tabular}{lcccc}
\toprule
\textbf{Backend} & \textbf{Watertight} & \textbf{Avg Verts} & \textbf{IoU} & \textbf{Speed} \\
\midrule
TripoSR (primary)       & 100\,\% & 92{,}311 & 0.452 & 5--18\,min \\
Visual Hull (fallback)  & 100\,\% & 49{,}991 & 0.21  & 0.5--2\,min \\
Sphere (debug)          & 100\,\% & $\sim$1{,}000 & --- & $<$1\,s \\
\bottomrule
\end{tabular}
\end{table}

\begin{figure}[htbp]
\centering
\begin{tikzpicture}
\begin{axis}[
  ybar, bar width=10pt,
  symbolic x coords={Watertight \%, Avg IoU (x100), Colour Err (inv.)},
  xtick=data,
  xticklabel style={font=\scriptsize, rotate=30, anchor=east},
  ymin=0, ymax=120,
  ylabel={Score (normalised to 100)},
  ylabel style={font=\scriptsize},
  nodes near coords,
  nodes near coords style={font=\tiny},
  height=5.0cm, width=\columnwidth,
  enlarge x limits=0.25,
  clip=false,
  legend style={font=\scriptsize, at={(0.98,0.98)}, anchor=north east},
  legend columns=2
]
\addplot[fill=blue!60] coordinates {
  ({Watertight \%},100)({Avg IoU (x100)},45)({Colour Err (inv.)},94)};
\addplot[fill=orange!70] coordinates {
  ({Watertight \%},100)({Avg IoU (x100)},21)({Colour Err (inv.)},75)};
\legend{TripoSR, Visual Hull}
\end{axis}
\end{tikzpicture}
\caption{Normalised quality comparison of TripoSR vs Visual Hull backends. Colour error is inverted
(higher bar = lower error = better). All backends achieve 100\,\% watertight output.}
\label{fig:backends}
\end{figure}

\subsection{Regression Test Suite}

The automated test suite comprises 82 passing tests (5 skipped due to GPU dependency) covering:
authentication flows, upload validation, processor module, reconstructor topology regressions,
mesh repair correctness, and end-to-end integration. All tests pass consistently on the development machine.

\begin{figure}[htbp]
\centering
\begin{tikzpicture}
\pie[
  radius=2.2,
  text=legend,
  color={blue!60, green!55, orange!60, red!50, purple!55},
  /tikz/every pin/.style={font=\scriptsize}
]{36/Auth \& Session (30),
  30/Processor module (25),
  23/Topology regression (19),
  16/Smoke \& integration (13),
  6/Skipped GPU (5)}
\end{tikzpicture}
\caption{Distribution of the 82-test regression suite by category.}
\label{fig:tests}
\end{figure}

% ─────────────────────────────────────────────────────────────────────────────
\subsection{Discussion}

\textbf{Colour accuracy.}
The region-based k-means approach achieves a mean colour error of 5.7 units (scale 0--255), which is perceptually
invisible to humans (the just-noticeable difference for colour in images is approximately 3--5 units). The primary
remaining artefact is that generated views are 20--30\,\% darker than the input photograph, which the per-view
calibration corrects for most objects. Objects with highly reflective surfaces (metallic knife blade) remain
challenging because specularity varies with viewpoint.

\textbf{Geometric fidelity.}
The silhouette IoU of 0.452 indicates that the reconstructed silhouette matches approximately 45\,\% of the
input-view silhouette area. This is consistent with the fundamental limitation of single-image reconstruction:
unseen surfaces must be hallucinated, and TripoSR's smooth-prior tendency produces slightly under-detailed
meshes for objects with complex, asymmetric geometry (thin structures, knife serrations).

\textbf{Processing time.}
The mean processing time of 575\,s (approximately 9.6 minutes) is dominated by TripoSR inference and the
six-view Zero123++ synthesis pass. On the development GPU (NVIDIA, CUDA 12.1), TripoSR inference takes
approximately 3--5 minutes. Faster inference is achievable with A100/H100 hardware or InstantMesh integration.

\textbf{Watertightness.}
The 100\,\% watertight rate on print exports is achieved by a four-stage repair pipeline: \texttt{fill\_holes},
\texttt{fix\_normals}, \texttt{fix\_winding}, and PyMeshFix. Any export failing the watertight check is rejected
before writing to disk.

% ─────────────────────────────────────────────────────────────────────────────
\section{Applications}

\begin{itemize}
  \item \textbf{E-commerce and product visualisation:} Automatic 3D models from a single product photograph,
        with lighting presets and screenshot tools for product-ready renders.
  \item \textbf{Gaming and VR:} Rapid asset prototyping from reference photographs, aided by the LOD0
        export path for real-time game engines.
  \item \textbf{Augmented reality:} GLB and AR/VR-optimised exports directly compatible with WebXR and
        ARCore for placing reconstructed real-world objects in digital environments.
  \item \textbf{Medical and anatomical modelling:} Conversion of photographs of anatomy, prosthetics,
        or surgical instruments into 3D models for training simulations.
  \item \textbf{3D printing:} The watertight print mesh and STL/3MF export paths enable rapid physical
        prototyping directly from a single photograph.
  \item \textbf{Robotics and simulation:} Simulation-oriented GLB and OBJ collision-mesh exports support
        importing reconstructed objects into simulators such as Gazebo and NVIDIA Isaac Sim.
\end{itemize}

% ─────────────────────────────────────────────────────────────────────────────
\section{Conclusion and Future Scope}

This paper presented TriGenix, an end-to-end, web-based single-image 3D reconstruction system developed and
evaluated across two iterative development phases. The Phase~1 pipeline combines Zero123++ multi-view diffusion
synthesis, Intel DPT-Large monocular depth estimation, and voxel-based Marching Cubes surface extraction to
produce watertight textured meshes. Phase~2 replaced the depth-fusion geometry stage with TripoSR neural
reconstruction and introduced a novel region-based k-means colour-projection algorithm.

Evaluated over 70 generated models, the system achieves a 100\,\% watertight rate on print exports, a mean
silhouette IoU of 0.452, and a mean colour error of 5.7 units — perceptually near-invisible. Nine export
formats cover printing, gaming, AR/VR, robotics, and texture-baking workflows. The system is production-grade,
with full multi-user authentication, real-time status polling, and a 82-test regression suite.

% ─────────────────────────────────────────────────────────────────────────────


\subsection{Limitations and Future Directions}%

The principal limitation is the hallucination problem common to all single-image reconstruction methods:
unseen surfaces are inferred rather than observed. Reconstruction quality degrades for objects with complex
topology, thin structures, or highly non-convex geometry.

Future work will:
\begin{enumerate}
  \item Integrate \textbf{InstantMesh}~\cite{instantmesh} to replace the TripoSR backend, targeting
        higher-fidelity meshes at sub-second generation time using FlexiCubes-based differentiable extraction.
  \item Add \textbf{Multi-Angle Capture Mode} (infrastructure already present as \texttt{multi\_angle\_processor.py})
        to accept four user-supplied photographs (front/back/left/right) and fuse them via posed reconstruction,
        substantially improving geometric accuracy.
  \item Implement \textbf{mesh quality scoring}, \textbf{wall thickness checking}, \textbf{auto-hollowing},
        \textbf{multi-part splitting}, and \textbf{multi-colour 3MF export} — all designed and ready for
        implementation in the next development phase.
  \item Extend the pipeline toward \textbf{near-real-time generation} via model quantisation and streaming
        Marching Cubes.
  \item Deploy on \textbf{cloud GPU infrastructure} (AWS, GCP) for wider public accessibility.
\end{enumerate}

% ─────────────────────────────────────────────────────────────────────────────

\section*{Acknowledgment}

The authors thank Prof. Sheetal Mhatre, Department of Data Science, Usha Mittal Institute of Technology,
S.N.D.T. Women's University, for continuous guidance and mentorship throughout the design, implementation,
and evaluation of TriGenix.

% ─────────────────────────────────────────────────────────────────────────────
\begin{thebibliography}{00}

\bibitem{survey_3d}
A. Kar, C. Häne, and J. Malik,
``Learning a Multi-View Stereo Machine,''
in \textit{Advances in Neural Information Processing Systems (NeurIPS)}, 2017.

\bibitem{mvs}
S. Seitz, B. Curless, J. Diebel, D. Scharstein, and R. Szeliski,
``A Comparison and Evaluation of Multi-View Stereo Reconstruction Algorithms,''
in \textit{Proc. IEEE/CVF CVPR}, vol.~1, pp.~519--528, 2006.

\bibitem{active_sensing}
B. Curless and M. Levoy,
``A Volumetric Method for Building Complex Models from Range Images,''
in \textit{Proc. ACM SIGGRAPH}, pp.~303--312, 1996.

\bibitem{zero123pp}
R. Shi, H. Chen, Z. Zhang, M. Liu, C. Xu, X. Wei, L. Chen, C. Zeng, and H. Su,
``Zero123++: A Single Image to Consistent Multi-view Diffusion Base Model,''
\textit{arXiv preprint arXiv:2310.15110}, 2023.

\bibitem{dpt}
R. Ranftl, A. Bochkovskiy, and V. Koltun, ``Vision Transformers for Dense Prediction,''
in \textit{Proc. IEEE/CVF International Conference on Computer Vision (ICCV)}, pp.~12179--12188, 2021.

\bibitem{lrm}
Y. Hong, K. Zhang, J. Gu, S. Bi, Y. Zhou, D. Liu, F. Liu, K. Sunkavalli, T. Bui, and H. Tan,
``LRM: Large Reconstruction Model for Single Image to 3D,''
in \textit{International Conference on Learning Representations (ICLR)}, 2024.

\bibitem{triposr}
D. Tochilkin, D. Pankratz, Z. Liu, Z. Huang, A. Letts, Y. Li, J.-X. Liang, C. Laforte, V. Jampani,
and Y. Cao,
``TripoSR: Fast 3D Object Reconstruction from a Single Image,''
\textit{arXiv preprint arXiv:2403.02156}, 2024.

\bibitem{instantmesh}
J. Xu, W. Cheng, Y. Gao, X. Wang, S. Gao, and Y. Shan,
``InstantMesh: Efficient 3D Mesh Generation from a Single Image with Sparse-view Large Reconstruction Models,''
\textit{arXiv preprint arXiv:2404.07191}, 2024.

\bibitem{eigen}
D. Eigen, C. Puhrsch, and R. Fergus, ``Depth Map Prediction from a Single Image using a Multi-Scale Deep Network,''
\textit{Advances in Neural Information Processing Systems (NeurIPS)}, vol.~27, pp.~2366--2374, 2014.

\bibitem{nerf}
B. Mildenhall, P.~P. Srinivasan, M. Tancik, J.~T. Barron, R. Ramamoorthi, and R. Ng,
``NeRF: Representing Scenes as Neural Radiance Fields for View Synthesis,''
in \textit{European Conference on Computer Vision (ECCV)}, pp.~405--421, Springer, 2020.

\bibitem{zero1to3}
R. Liu, R. Wu, B. Van Hoorick, P. Tokmakov, S. Zakharov, and C. Vondrick,
``Zero-1-to-3: Zero-shot One Image to 3D Object,''
in \textit{Proc. IEEE/CVF ICCV}, pp.~9298--9309, 2023.

\bibitem{marching_cubes}
W.~E. Lorensen and H.~E. Cline,
``Marching Cubes: A High Resolution 3D Surface Construction Algorithm,''
\textit{ACM SIGGRAPH Computer Graphics}, vol.~21, no.~4, pp.~163--169, 1987.

\end{thebibliography}

\end{document}
